\documentclass[UTF8]{ctexart}
\usepackage{amsmath}
\usepackage{tocloft}
\usepackage[bigfiles]{pdfbase}
\usepackage{amsthm,amssymb}
\usepackage{booktabs}
\usepackage{graphicx}
\usepackage[hidelinks]{hyperref}
\usepackage{geometry}
\usepackage{float}
\usepackage{multirow}
\usepackage{indentfirst}
\usepackage{diagbox}
\usepackage{listings}
\usepackage{xcolor}
\usepackage{markdown}
\definecolor{codegreen}{rgb}{0,0.6,0}
\definecolor{codegray}{rgb}{0.5,0.5,0.5}
\definecolor{codepurple}{rgb}{0.58,0,0.82}
\definecolor{backcolour}{rgb}{0.95,0.95,0.92}

\lstdefinestyle{mystyle}{
    backgroundcolor=\color{backcolour},   
    commentstyle=\color{codegreen},
    keywordstyle=\color{magenta},
    numberstyle=\tiny\color{codegray},
    stringstyle=\color{codepurple},
    basicstyle=\ttfamily\footnotesize,
    breakatwhitespace=false,         
    breaklines=true,                 
    captionpos=b,                    
    keepspaces=true,                 
    numbers=left,                    
    numbersep=5pt,                  
    showspaces=false,                
    showstringspaces=false,
    showtabs=false,                  
    tabsize=2
}
\lstset{style=mystyle}
\newcommand\question[2]{\noindent\fbox{\parbox[t]{\linewidth}{\hspace{0pt}{\textbf{#1\quad}#2}}}}


\usepackage{graphicx}
\usepackage{setspace}
\renewcommand{\cftsecleader}{\cftdotfill{\cftdotsep}}
\geometry{a4paper,left=3cm,right=3cm,top=2cm,bottom=2cm} 


\CTEXsetup[format={\Large \bfseries}]{section}
\title{\textbf{人工智能赋能的数字集成电路设计 作业一}}
\author{陈天乐\ \ [student ID removed] }
\date{\today}
\begin{document}
\maketitle
\tableofcontents
\newpage
\section{One Shot Prompt}

我在原本的prompt基础上添加了几个Tips，并且加入了Simulation Result， Last RTL， Last TestBench，分别是上一次代码的仿真结果、生成的rtl代码
和testbench代码。

\begin{lstlisting}[frame=shadowbox]
Below is a template for the problem input and response. You should reference this response format.

---

Here is an examplaray prompt input:
Please help me generate a FIFO, including RTL code and TestBench.

## Problem Description

ModelName: FIFO

InputNode:

- clk: Clock signal (rising edge effective).
- rst_n: Active-low reset signal.
- data_in: 32-bit input data.
- wr_en: Write enable signal, high effective.
- rd_en: Read enable signal, high effective.

OutputNode:

- data_out: 32-bit output data.
- full: FIFO full signal, high effective.
- empty: FIFO empty signal, high effective.

Specification:

- FIFO Size: 16 entries, each entry is 32 bits wide.
- Behavior:
  - When wr_en is high and full is low, the input data data_in is written into the FIFO.
  - When rd_en is high and empty is low, the oldest data in the FIFO is read out as data_out.
  - If wr_en is high and full is high, the write operation is ignored.
  - If rd_en is high and empty is high, the data_out remains unchanged.
- TestBench should contain the following test conditions:
  - Basic Write and Read Operations: Write a few data entries and read them back in sequence.
  - Concurrent Write and Read: Write and read data simultaneously to test FIFO behavior.
  - FIFO Full Condition: Write data continuously until the FIFO is full, then attempt to write more data.
  - FIFO Empty Condition: Read data continuously until the FIFO is empty, then attempt to read more data.
  - Mixed Write and Read Patterns: Write and read data in various patterns to test FIFO robustness.
  - Edge Cases: Write maximum number of entries and then read them back; read all entries and then write again.
  - Reset Behavior: Test FIFO behavior after a reset signal is asserted.


Here is an examplaray prompt output:

### RTL

```verilog
module FIFO (
    input wire clk,
    input wire rst_n,
    input wire [31:0] data_in,
    input wire wr_en,
    input wire rd_en,
    output reg [31:0] data_out,
    output wire full,
    output wire empty
);

    // FIFO parameters
    parameter DEPTH = 16;
    parameter WIDTH = 32;
    
    // Internal memory
    reg [WIDTH-1:0] memory [0:DEPTH-1];
    
    // Pointers
    reg [3:0] wr_ptr;
    reg [3:0] rd_ptr;
    reg [4:0] count;  // Need 5 bits to count 0-16
    
    // Full and empty flags
    assign full = (count == DEPTH);
    assign empty = (count == 0);
    
    // FIFO control logic
    always @(posedge clk or negedge rst_n) begin
        if (!rst_n) begin
            // Reset all pointers and memory
            wr_ptr <= 0;
            rd_ptr <= 0;
            count <= 0;
            data_out <= 0;
            
            // Optional: Clear memory on reset
            for (integer i = 0; i < DEPTH; i = i + 1) begin
                memory[i] <= 0;
            end
        end else begin
            // Write operation
            if (wr_en && !full) begin
                memory[wr_ptr] <= data_in;
                wr_ptr <= wr_ptr + 1;
                if (wr_ptr == DEPTH-1) wr_ptr <= 0;
            end
            
            // Read operation
            if (rd_en && !empty) begin
                data_out <= memory[rd_ptr];
                rd_ptr <= rd_ptr + 1;
                if (rd_ptr == DEPTH-1) rd_ptr <= 0;
            end
            
            // Update count
            case ({wr_en && !full, rd_en && !empty})
                2'b00: count <= count;  // No operation
                2'b01: count <= count - 1;  // Read only
                2'b10: count <= count + 1;  // Write only
                2'b11: count <= count;  // Simultaneous read and write
            endcase
        end
    end

endmodule
```

### TestBench

```systemverilog
module FIFO_tb;
    // Parameters
    parameter CLK_PERIOD = 10;
    parameter DEPTH = 16;
    
    // Signals
    reg clk;
    reg rst_n;
    reg [31:0] data_in;
    reg wr_en;
    reg rd_en;
    wire [31:0] data_out;
    wire full;
    wire empty;
    
    // Test variables
    reg [31:0] expected_data;
    integer error_count;
    integer test_case;
    
    // Instantiate DUT
    FIFO dut (
        .clk(clk),
        .rst_n(rst_n),
        .data_in(data_in),
        .wr_en(wr_en),
        .rd_en(rd_en),
        .data_out(data_out),
        .full(full),
        .empty(empty)
    );
    
    // Clock generation
    always #(CLK_PERIOD/2) clk = ~clk;
    
    // Waveform dumping
    initial begin
        $dumpfile("Out/wave.vcd");
        $dumpvars(0, FIFO_tb);
    end
    
    // Test tasks
    task write_data(input [31:0] data);
        begin
            data_in = data;
            wr_en = 1;
            @(posedge clk);
            wr_en = 0;
        end
    endtask
    
    task read_data(output [31:0] data);
        begin
            rd_en = 1;
            @(posedge clk);
            rd_en = 0;
            data = data_out;
        end
    endtask
    
    task reset_fifo;
        begin
            rst_n = 0;
            wr_en = 0;
            rd_en = 0;
            data_in = 0;
            @(posedge clk);
            @(posedge clk);
            rst_n = 1;
            @(posedge clk);
        end
    endtask
    
    task check_result(input [31:0] actual, input [31:0] expected, input string test_name);
        begin
            if (actual !== expected) begin
                $display("ERROR: %s - Expected: %h, Got: %h", test_name, expected, actual);
                error_count = error_count + 1;
            end
        end
    endtask
    
    // Main test sequence
    initial begin
        // Initialize
        clk = 0;
        error_count = 0;
        test_case = 0;
        
        // Test 1: Reset behavior
        test_case = 1;
        $display("Test %0d: Reset Behavior", test_case);
        reset_fifo();
        check_result(data_out, 32'h0, "Reset data_out");
        check_result(empty, 1'b1, "Reset empty flag");
        check_result(full, 1'b0, "Reset full flag");
        
        // Test 2: Basic write and read operations
        test_case = 2;
        $display("Test %0d: Basic Write and Read Operations", test_case);
        reset_fifo();
        
        // Write 4 data items
        write_data(32'h12345678);
        write_data(32'hABCDEF01);
        write_data(32'hDEADBEEF);
        write_data(32'hCAFEBABE);
        
        // Read them back
        read_data(expected_data);
        check_result(expected_data, 32'h12345678, "First read");
        read_data(expected_data);
        check_result(expected_data, 32'hABCDEF01, "Second read");
        read_data(expected_data);
        check_result(expected_data, 32'hDEADBEEF, "Third read");
        read_data(expected_data);
        check_result(expected_data, 32'hCAFEBABE, "Fourth read");
        
        // Test 3: Concurrent write and read
        test_case = 3;
        $display("Test %0d: Concurrent Write and Read", test_case);
        reset_fifo();
        
        fork
            begin
                for (integer i = 0; i < 8; i = i + 1) begin
                    write_data(32'h1000 + i);
                end
            end
            begin
                for (integer i = 0; i < 8; i = i + 1) begin
                    read_data(expected_data);
                    check_result(expected_data, 32'h1000 + i, $sformatf("Concurrent read %0d", i));
                end
            end
        join
        
        // Test 4: FIFO full condition
        test_case = 4;
        $display("Test %0d: FIFO Full Condition", test_case);
        reset_fifo();
        
        // Fill the FIFO
        for (integer i = 0; i < DEPTH; i = i + 1) begin
            write_data(32'h2000 + i);
            if (i == DEPTH-1) begin
                check_result(full, 1'b1, "FIFO should be full");
            end
        end
        
        // Try to write when full (should be ignored)
        data_in = 32'hBAD0BAD0;
        wr_en = 1;
        @(posedge clk);
        wr_en = 0;
        @(posedge clk);
        check_result(full, 1'b1, "FIFO should remain full");
        
        // Test 5: FIFO empty condition
        test_case = 5;
        $display("Test %0d: FIFO Empty Condition", test_case);
        
        // Empty the FIFO
        for (integer i = 0; i < DEPTH; i = i + 1) begin
            read_data(expected_data);
            check_result(expected_data, 32'h2000 + i, $sformatf("Emptying read %0d", i));
            if (i == DEPTH-1) begin
                check_result(empty, 1'b1, "FIFO should be empty");
            end
        end
        
        // Try to read when empty
        rd_en = 1;
        @(posedge clk);
        rd_en = 0;
        @(posedge clk);
        check_result(empty, 1'b1, "FIFO should remain empty");
        
        // Test 6: Mixed write and read patterns
        test_case = 6;
        $display("Test %0d: Mixed Write and Read Patterns", test_case);
        reset_fifo();
        
        // Pattern 1: Write 3, read 2, write 2, read 3
        write_data(32'h3001); write_data(32'h3002); write_data(32'h3003);
        read_data(expected_data); check_result(expected_data, 32'h3001, "Mixed pattern 1a");
        read_data(expected_data); check_result(expected_data, 32'h3002, "Mixed pattern 1b");
        write_data(32'h3004); write_data(32'h3005);
        read_data(expected_data); check_result(expected_data, 32'h3003, "Mixed pattern 1c");
        read_data(expected_data); check_result(expected_data, 32'h3004, "Mixed pattern 1d");
        read_data(expected_data); check_result(expected_data, 32'h3005, "Mixed pattern 1e");
        
        // Test 7: Edge cases
        test_case = 7;
        $display("Test %0d: Edge Cases", test_case);
        reset_fifo();
        
        // Write maximum entries and read them back
        for (integer i = 0; i < DEPTH; i = i + 1) begin
            write_data(32'h4000 + i);
        end
        check_result(full, 1'b1, "Edge case - full");
        
        for (integer i = 0; i < DEPTH; i = i + 1) begin
            read_data(expected_data);
            check_result(expected_data, 32'h4000 + i, $sformatf("Edge case read %0d", i));
        end
        check_result(empty, 1'b1, "Edge case - empty");
        
        // Write again after reading all
        write_data(32'h5001);
        read_data(expected_data);
        check_result(expected_data, 32'h5001, "Write after complete read");
        
        // Final result
        if (error_count == 0) begin
            $display("TestSuccessed: All tests passed!");
        end else begin
            $display("TestFailed: %0d errors found", error_count);
        end
        
        $finish;
    end

endmodule
```

---

Here is the real input prompt

Please help me generate [ModelName], including RTL code and TestBench.

## Problem Description

ModelName: [ModelName]

InputNode: [InputNode]

OutputNode: [OutputNode]

Specification: [Specification]

TestBench: [TestSpec]

Last_RTL: [LastRTL]

Last_TestBench: [LastTestBench]

SimulationResult: [SimulationResult]


### Tips

Here is some tips that you need to follow:

1. Use Verilog for RTL and SystemVerilog for TestBench.
2. We use the -g2005sv specification supported by iverilog for compilation. Please make sure that the verilog used can be supported. Please ensure that the system functions are supported by the -g2005sv. Some functions like $sformatf() will arise error.
3. Unpacked structs not supported in iverilog, do not use typedef struct.
4. When using the $bitstoreal function, please make sure that the parameter is 64bit.
5. Save the simulated waveform diagram with $dumpfile and $dumpvars instructions. Save the vcd file as Out/wave.vcd
6. In Testbench, please verify whether the function is correct. If you pass all the tests, please output TestSuccessed, otherwise you will output TestFailed.
7. In Verilog, for loops cannot be used to initialize arrays within always blocks for sequential logic that are triggered by a clock.
8. In Verilog, earlier versions (such as those prior to Verilog-2001) do not support directly declaring loop variables in the initialization section of a for loop. This is an vital part, so you should pay attention to this point. Otherwise the simulator will call error: name is not a valid net. You can define it previously.
9. Last_RTL, Last_TestBench, and SimulationResult refer to the previously generated RTL code, TestBench code, and the error results from the simulator. Please adjust the newly generated code by incorporating these, and make every effort to eliminate the errors.
10. When executing the Testbench, first print out the line: "### SimulationResult\n '''markdown\n" in testbench.


\end{lstlisting}

\section{Async FIFO Specification}


\begin{figure}[H]
	\centering
	\includegraphics[width=\textwidth]{spec1.png}
\end{figure}


\begin{figure}[H]
	\centering
	\includegraphics[width=\textwidth]{spec2.png}
\end{figure}


\begin{figure}[H]
	\centering
	\includegraphics[width=\textwidth]{spec3.png}
\end{figure}


\begin{figure}[H]
	\centering
	\includegraphics[width=\textwidth]{spec4.png}
\end{figure}


\begin{figure}[H]
	\centering
	\includegraphics[width=\textwidth]{spec5.png}
\end{figure}

\section{RTL and TestBench}

\begin{lstlisting}[frame=shadowbox,language=verilog,caption=async\_fifo.v]
module async_fifo (
    input wire wr_clk,
    input wire rd_clk,
    input wire rst_n,
    input wire wr_en,
    input wire rd_en,
    input wire [15:0] data_in,
    output reg [15:0] data_out,
    output wire full,
    output wire empty,
    output wire almost_full,
    output wire almost_empty
);

    // Parameters
    parameter DEPTH = 16;
    parameter WIDTH = 16;
    parameter ALMOST_FULL_THRESHOLD = 14;
    parameter ALMOST_EMPTY_THRESHOLD = 2;
    
    // Internal memory
    reg [WIDTH-1:0] memory [0:DEPTH-1];
    
    // Pointers (5 bits for depth 16)
    reg [4:0] wr_ptr_bin;
    reg [4:0] rd_ptr_bin;
    wire [4:0] wr_ptr_gray;
    wire [4:0] rd_ptr_gray;
    
    // Synchronized pointers
    reg [4:0] wr_ptr_gray_sync1, wr_ptr_gray_sync2;
    reg [4:0] rd_ptr_gray_sync1, rd_ptr_gray_sync2;
    
    // Binary synchronized pointers
    reg [4:0] wr_ptr_bin_sync;
    reg [4:0] rd_ptr_bin_sync;
    
    // FIFO counts
    reg [4:0] fifo_count_wr_domain;
    reg [4:0] fifo_count_rd_domain;
    
    // Convert binary to Gray code
    assign wr_ptr_gray = wr_ptr_bin ^ (wr_ptr_bin >> 1);
    assign rd_ptr_gray = rd_ptr_bin ^ (rd_ptr_bin >> 1);
    
    // Convert Gray code to binary
    function [4:0] gray2bin;
        input [4:0] gray;
        begin
            gray2bin[4] = gray[4];
            gray2bin[3] = gray[4] ^ gray[3];
            gray2bin[2] = gray[4] ^ gray[3] ^ gray[2];
            gray2bin[1] = gray[4] ^ gray[3] ^ gray[2] ^ gray[1];
            gray2bin[0] = gray[4] ^ gray[3] ^ gray[2] ^ gray[1] ^ gray[0];
        end
    endfunction
    
    // Write pointer synchronization (to read clock domain)
    always @(posedge rd_clk or negedge rst_n) begin
        if (!rst_n) begin
            wr_ptr_gray_sync1 <= 0;
            wr_ptr_gray_sync2 <= 0;
            wr_ptr_bin_sync <= 0;
        end else begin
            wr_ptr_gray_sync1 <= wr_ptr_gray;
            wr_ptr_gray_sync2 <= wr_ptr_gray_sync1;
            wr_ptr_bin_sync <= gray2bin(wr_ptr_gray_sync2);
        end
    end
    
    // Read pointer synchronization (to write clock domain)
    always @(posedge wr_clk or negedge rst_n) begin
        if (!rst_n) begin
            rd_ptr_gray_sync1 <= 0;
            rd_ptr_gray_sync2 <= 0;
            rd_ptr_bin_sync <= 0;
        end else begin
            rd_ptr_gray_sync1 <= rd_ptr_gray;
            rd_ptr_gray_sync2 <= rd_ptr_gray_sync1;
            rd_ptr_bin_sync <= gray2bin(rd_ptr_gray_sync2);
        end
    end
    
    // Calculate FIFO count in write domain
    always @(*) begin
        if (wr_ptr_bin[4] == rd_ptr_bin_sync[4]) begin
            fifo_count_wr_domain = wr_ptr_bin[3:0] - rd_ptr_bin_sync[3:0];
        end else begin
            fifo_count_wr_domain = DEPTH - rd_ptr_bin_sync[3:0] + wr_ptr_bin[3:0];
        end
    end
    
    // Calculate FIFO count in read domain
    always @(*) begin
        if (wr_ptr_bin_sync[4] == rd_ptr_bin[4]) begin
            fifo_count_rd_domain = wr_ptr_bin_sync[3:0] - rd_ptr_bin[3:0];
        end else begin
            fifo_count_rd_domain = DEPTH - rd_ptr_bin[3:0] + wr_ptr_bin_sync[3:0];
        end
    end
    
    // Generate flags
    assign full = (fifo_count_wr_domain == DEPTH);
    assign empty = (fifo_count_rd_domain == 0);
    assign almost_full = (fifo_count_wr_domain >= ALMOST_FULL_THRESHOLD);
    assign almost_empty = (fifo_count_rd_domain <= ALMOST_EMPTY_THRESHOLD);
    
    // Write control logic
    always @(posedge wr_clk or negedge rst_n) begin
        integer i;
        if (!rst_n) begin
            wr_ptr_bin <= 0;
        end else begin
            if (wr_en && !full) begin
                memory[wr_ptr_bin[3:0]] <= data_in;
                wr_ptr_bin <= wr_ptr_bin + 1;
            end
        end
    end
    
    // Read control logic
    always @(posedge rd_clk or negedge rst_n) begin
        if (!rst_n) begin
            rd_ptr_bin <= 0;
            data_out <= 0;
        end else begin
            if (rd_en && !empty) begin
                data_out <= memory[rd_ptr_bin[3:0]];
                rd_ptr_bin <= rd_ptr_bin + 1;
            end
        end
    end
    
    // Initialize memory (for simulation only)
    integer init_i;
    initial begin
        for (init_i = 0; init_i < DEPTH; init_i = init_i + 1) begin
            memory[init_i] = 0;
        end
    end

endmodule
\end{lstlisting}

\begin{lstlisting}[frame=shadowbox, language=verilog, caption=tb\_async\_fifo.sv]
module async_fifo_tb;
    // Parameters
    parameter WR_CLK_PERIOD = 10; // 100MHz
    parameter RD_CLK_PERIOD = 20; // 50MHz
    parameter DEPTH = 16;
    
    // Signals
    reg wr_clk;
    reg rd_clk;
    reg rst_n;
    reg wr_en;
    reg rd_en;
    reg [15:0] data_in;
    wire [15:0] data_out;
    wire full;
    wire empty;
    wire almost_full;
    wire almost_empty;
    
    // Test variables
    reg [15:0] expected_data;
    integer error_count;
    integer test_case;
    integer i;
    
    // Instantiate DUT
    async_fifo dut (
        .wr_clk(wr_clk),
        .rd_clk(rd_clk),
        .rst_n(rst_n),
        .wr_en(wr_en),
        .rd_en(rd_en),
        .data_in(data_in),
        .data_out(data_out),
        .full(full),
        .empty(empty),
        .almost_full(almost_full),
        .almost_empty(almost_empty)
    );
    
    // Clock generation
    always #(WR_CLK_PERIOD/2) wr_clk = ~wr_clk;
    always #(RD_CLK_PERIOD/2) rd_clk = ~rd_clk;
    
    // Waveform dumping
    initial begin
        $dumpfile("Out/wave.vcd");
        $dumpvars(0, async_fifo_tb);
    end
    
    // Test tasks
    task write_data(input [15:0] data);
        begin
            data_in = data;
            wr_en = 1;
            @(negedge wr_clk); // Ensure wr_en setup time
            @(posedge wr_clk);
            wr_en = 0;
        end
    endtask
    
    task read_data(output [15:0] data);
        begin
            rd_en = 1;
            @(posedge rd_clk);
            rd_en = 0;
            data = data_out;
        end
    endtask
    
    task reset_fifo;
        begin
            rst_n = 0;
            wr_en = 0;
            rd_en = 0;
            data_in = 0;
            @(posedge wr_clk);
            @(posedge rd_clk);
            @(posedge wr_clk);
            @(posedge rd_clk);
            rst_n = 1;
            @(posedge wr_clk);
            @(posedge rd_clk);
        end
    endtask
    
    task check_result(input actual, input expected, input string test_name);
        begin
            if (actual !== expected) begin
                $display("ERROR: %s - Expected: %b, Got: %b", test_name, expected, actual);
                error_count = error_count + 1;
            end
        end
    endtask
    
    task check_data_result(input [15:0] actual, input [15:0] expected, input string test_name);
        begin
            if (actual !== expected) begin
                $display("ERROR: %s - Expected: %h, Got: %h", test_name, expected, actual);
                error_count = error_count + 1;
            end
        end
    endtask
    
    // Main test sequence
    initial begin
        // Initialize
        wr_clk = 0;
        rd_clk = 0;
        error_count = 0;
        test_case = 0;
        
        $display("### SimulationResult");
        $display("'''markdown");
        
        // Test 1: Reset behavior
        test_case = 1;
        $display("Test %0d: Reset Behavior", test_case);
        reset_fifo();
        check_data_result(data_out, 16'h0, "Reset data_out");
        check_result(empty, 1'b1, "Reset empty flag");
        check_result(full, 1'b0, "Reset full flag");
        check_result(almost_empty, 1'b1, "Reset almost_empty flag");
        check_result(almost_full, 1'b0, "Reset almost_full flag");
        if (error_count == 0) begin
            $display("Test %0d: Reset Behavior - PASSED", test_case);
        end
        
        // Test 2: Basic write and read operations
        test_case = 2;
        $display("Test %0d: Basic Write and Read Operations", test_case);
        reset_fifo();
        
        // Write 4 data items
        write_data(16'h1234);
        write_data(16'hABCD);
        write_data(16'h5678);
        write_data(16'hEF01);
        
        // Add sufficient delay for synchronization (more cycles for async FIFO)
        repeat(8) @(posedge rd_clk);
        
        // Read them back
        read_data(expected_data);
        check_data_result(expected_data, 16'h1234, "First read");
        read_data(expected_data);
        check_data_result(expected_data, 16'hABCD, "Second read");
        read_data(expected_data);
        check_data_result(expected_data, 16'h5678, "Third read");
        read_data(expected_data);
        check_data_result(expected_data, 16'hEF01, "Fourth read");
        if (error_count == 0) begin
            $display("Test %0d: Basic Write and Read Operations - PASSED", test_case);
        end
        
        // Test 3: Asynchronous clock test
        test_case = 3;
        $display("Test %0d: Asynchronous Clock Test", test_case);
        reset_fifo();
        
        fork
            begin
                for (i = 0; i < 8; i = i + 1) begin
                    write_data(16'h1000 + i);
                end
            end
            begin
                // Add more delay to ensure first data is written and synchronized
                repeat(10) @(posedge rd_clk);
                for (i = 0; i < 8; i = i + 1) begin
                    read_data(expected_data);
                    check_data_result(expected_data, 16'h1000 + i, "Async read");
                end
            end
        join
        if (error_count == 0) begin
            $display("Test %0d: Asynchronous Clock Test - PASSED", test_case);
        end
        
        // Test 4: FIFO full condition
        test_case = 4;
        $display("Test %0d: FIFO Full Condition", test_case);
        reset_fifo();
        
        // Fill the FIFO
        for (i = 0; i < DEPTH; i = i + 1) begin
            write_data(16'h2000 + i);
            if (i == DEPTH-1) begin
                check_result(full, 1'b1, "FIFO should be full");
                check_result(almost_full, 1'b1, "FIFO should be almost full");
            end
        end
        
        // Try to write when full (should be ignored)
        data_in = 16'hBAD0;
        wr_en = 1;
        @(negedge wr_clk);
        @(posedge wr_clk);
        wr_en = 0;
        @(posedge wr_clk);
        check_result(full, 1'b1, "FIFO should remain full");
        if (error_count == 0) begin
            $display("Test %0d: FIFO Full Condition - PASSED", test_case);
        end
        
        // Test 5: FIFO empty condition
        test_case = 5;
        $display("Test %0d: FIFO Empty Condition", test_case);
        reset_fifo();
        
        // Fill the FIFO first
        for (i = 0; i < DEPTH; i = i + 1) begin
            write_data(16'h2000 + i);
        end
        
        // Add delay for synchronization
        repeat(8) @(posedge rd_clk);
        
        // Empty the FIFO
        for (i = 0; i < DEPTH; i = i + 1) begin
            read_data(expected_data);
            check_data_result(expected_data, 16'h2000 + i, "Emptying read");
            if (i == DEPTH-1) begin
                check_result(empty, 1'b1, "FIFO should be empty");
                check_result(almost_empty, 1'b1, "FIFO should be almost empty");
            end
        end
        
        // Try to read when empty
        rd_en = 1;
        @(posedge rd_clk);
        rd_en = 0;
        @(posedge rd_clk);
        check_result(empty, 1'b1, "FIFO should remain empty");
        if (error_count == 0) begin
            $display("Test %0d: FIFO Empty Condition - PASSED", test_case);
        end
        
        // Test 6: Almost-full/almost-empty threshold test
        test_case = 6;
        $display("Test %0d: Almost-Full/Almost-Empty Threshold Test", test_case);
        reset_fifo();
        
        // Write until almost full threshold (14)
        for (i = 0; i < 13; i = i + 1) begin
            write_data(16'h3000 + i);
        end
        check_result(almost_full, 1'b0, "Should not be almost full yet");
        write_data(16'h3000 + 13);
        check_result(almost_full, 1'b1, "Should be almost full at 14 entries");
        check_result(full, 1'b0, "Should not be full yet");
        
        // Add delay for synchronization
        repeat(8) @(posedge rd_clk);
        
        // Read until almost empty threshold (2)
        for (i = 0; i < 12; i = i + 1) begin
            read_data(expected_data);
            check_data_result(expected_data, 16'h3000 + i, "Threshold read");
        end
        // After reading 12 from 14, we have 2 left - should be almost empty
        check_result(almost_empty, 1'b1, "Should be almost empty at 2 entries");
        check_result(empty, 1'b0, "Should not be empty yet");
        
        read_data(expected_data);
        check_result(almost_empty, 1'b1, "Should still be almost empty at 1 entry");
        check_result(empty, 1'b0, "Should not be empty yet");
        
        if (error_count == 0) begin
            $display("Test %0d: Almost-Full/Almost-Empty Threshold Test - PASSED", test_case);
        end
        
        // Test 7: Concurrent read-write test
        test_case = 7;
        $display("Test %0d: Concurrent Read-Write Test", test_case);
        reset_fifo();
        
        fork
            begin
                for (i = 0; i < 10; i = i + 1) begin
                    write_data(16'h4000 + i);
                end
            end
            begin
                // Add more delay to ensure first data is written and synchronized
                repeat(12) @(posedge rd_clk);
                for (i = 0; i < 10; i = i + 1) begin
                    read_data(expected_data);
                    check_data_result(expected_data, 16'h4000 + i, "Concurrent read");
                end
            end
        join
        if (error_count == 0) begin
            $display("Test %0d: Concurrent Read-Write Test - PASSED", test_case);
        end
        
        // Test 8: Edge case test
        test_case = 8;
        $display("Test %0d: Edge Case Test", test_case);
        reset_fifo();
        
        // Write 15 entries (one before full)
        for (i = 0; i < 15; i = i + 1) begin
            write_data(16'h5000 + i);
        end
        check_result(full, 1'b0, "Should not be full at 15 entries");
        check_result(almost_full, 1'b1, "Should be almost full at 15 entries");
        
        // Write 16th entry (full)
        write_data(16'h5000 + 15);
        check_result(full, 1'b1, "Should be full at 16 entries");
        
        // Add delay for synchronization
        repeat(8) @(posedge rd_clk);
        
        // Read 14 entries - should have 2 left (almost empty)
        for (i = 0; i < 14; i = i + 1) begin
            read_data(expected_data);
            check_data_result(expected_data, 16'h5000 + i, "Edge case read");
        end
        check_result(almost_empty, 1'b1, "Should be almost empty at 2 entries");
        
        // Read 1 more entry - should have 1 left (still almost empty)
        read_data(expected_data);
        check_result(almost_empty, 1'b1, "Should still be almost empty at 1 entry");
        
        // Read last entry - should be empty
        read_data(expected_data);
        check_result(empty, 1'b1, "Should be empty after reading all");
        check_result(almost_empty, 1'b1, "Should be almost empty at 0 entries");
        
        if (error_count == 0) begin
            $display("Test %0d: Edge Case Test - PASSED", test_case);
        end
        
        // Final result
        if (error_count == 0) begin
            $display("TestSuccessed: All tests passed!");
        end else begin
            $display("TestFailed: %0d errors found", error_count);
        end
        
        $display("'''");
        $finish;
    end

endmodule
\end{lstlisting}

\section{Simulation Result}
一开始没有将上次生成的代码和仿真报错结果写进prompt里，导致模型会一直重复在几个错误版本之间反复跳跃，而不会有实质性的改进。在让模型有记忆性后，得到的仿真结果如下：

\begin{lstlisting}[frame=shadowbox]
Test 1: Reset Behavior
Test 1: Reset Behavior - PASSED
Test 2: Basic write and Read operations 
Test 2: Basic write and Read operations - PASSED
Test 3: Asynchronous clock Test
Test 3: Asynchronous clock Test - PASSED
Test 4: FIF0 Full condition
Test 4: FIF0 Full condition - PASSED
Test 5: FIFO Empty condition
Test S: FIFO Empty condition -PASSED
Test 6: Almost-Full/Almost-Empty Threshold Test
Test 6: Almost-Full/Almost-Empty Threshold Test - PASSED
Test 7: Concurrent Read-write Test
Test 7:concurrent Read-write Test - PASSED
Test 8: Edge case Test
Test 8: Edge case Test - PASSED
Testsuccessed: All tests passed!
\end{lstlisting}

现在的模型可以一次性生成可以运行且通过了所有测试点的异步FIFO代码了。仿真波形图如下：

\begin{figure}[H]
	\centering
	\includegraphics[width=\textwidth]{sim.png}
	\caption{仿真结果和波形图}
\end{figure}

\end{document}